% Bloom levels: remember, understand, apply, analyze, evaluate, create.
\begin{problema}\label{prob:1d85679}
Let $X_1,X_2,\ldots,X_n$ be an independent random sample from a population whose random variable satisfies $E[X^2]=\theta$, where $\theta>0$ is the population second raw moment. Consider
\[
	\hat\theta=\frac1n\sum_{i=1}^n X_i^2.
\]
Prove analytically, using linearity of expectation, that $\hat\theta$ is an unbiased estimator of $\theta$.
\end{problema}

\begin{problema}\label{prob:6f639cf}
The time between failures of cloud data-center servers is modeled by a continuous random variable $X$ with Exponential distribution of parameter $\lambda>0$:
\[
	f(x\mid\lambda)=\lambda e^{-\lambda x},\qquad x>0.
\]
\begin{enumerate}[label=(\alph*)]
	\item Obtain the method-of-moments estimator $\hat\lambda_{MoM}$.
	\item Obtain the maximum-likelihood estimator $\hat\lambda_{MLE}$. Do both methods coincide in this model? Explain a general condition under which method of moments and maximum likelihood coincide.
\end{enumerate}
\end{problema}

\begin{problema}\label{prob:0363afd}
In a social-media advertising campaign, whether the $i$th user clicks an ad is modeled as $X_i\sim\operatorname{Bernoulli}(p)$, where $p\in(0,1)$ is the unknown true click-through rate and $n$ is the number of independent impressions.
\begin{enumerate}[label=(\alph*)]
	\item Derive the maximum-likelihood estimator $\hat p_{MLE}$ from the likelihood of binary trials.
	\item Compute the exact variance $\Var(\hat p_{MLE})$ in terms of $p$ and $n$, and state the estimated standard error $\operatorname{SE}(\hat p_{MLE})$.
\end{enumerate}
\end{problema}

\begin{problema}\label{prob:977b8dd}
\textbf{Absolute efficiency and the Cramér--Rao lower bound:} let $X_1,\ldots,X_n$ be a sample from a Poisson distribution with rate $\lambda>0$.
\begin{enumerate}[label=(\alph*)]
	\item Compute the Fisher information $I(\lambda)$ in one observation $X\sim\operatorname{Poisson}(\lambda)$ by differentiating the log-density and taking the expected second derivative.
	\item Determine the Cramér--Rao lower bound for the variance of any unbiased estimator of $\lambda$ based on a sample of size $n$.
	\item Prove that $\hat\lambda_{MLE}=\bar X$ attains the bound exactly, and conclude that it is the uniformly minimum-variance unbiased estimator.
\end{enumerate}
\end{problema}

\begin{problema}\label{prob:23bb182}
Let $X_1,\ldots,X_n$ be a random sample from a Uniform distribution on $[0,\theta]$, where $\theta>0$ is unknown.
\begin{enumerate}[label=(\alph*)]
	\item Find the method-of-moments estimator $\hat\theta_{MoM}$ and the maximum-likelihood estimator $\hat\theta_{MLE}=X_{(n)}=\max(X_1,\ldots,X_n)$. Briefly justify why the usual procedure of differentiating and setting the derivative equal to zero fails here.
	\item Compare the biases of $\hat\theta_{MoM}$ and $\hat\theta_{MLE}$, and construct from the latter a modified estimator $\hat\theta^*$ that is unbiased. Judge which of the three estimators you would recommend in practice and why.
\end{enumerate}
\end{problema}

\begin{problema}\label{prob:0675127}
Design an original problem that illustrates the link between classical point estimation and training models in data science: for a single binary observation $y\in\{0,1\}$ with predicted probability $p$, prove that minimizing the cross-entropy loss
\[
	\mathcal L(p)=-[y\ln p+(1-y)\ln(1-p)]
\]
with respect to $p$ is mathematically equivalent to maximizing the log-likelihood of a Bernoulli model, and therefore produces the same estimator as classical maximum likelihood.
\end{problema}

\begin{solproblema}[prob:1d85679]
By linearity of expectation,
\[
	E[\hat\theta]
	=E\left[\frac1n\sum_{i=1}^n X_i^2\right]
	=\frac1n\sum_{i=1}^n E[X_i^2].
\]
Since the observations are iid and each has second raw moment $\theta$,
\[
	E[\hat\theta]=\frac1n(n\theta)=\theta.
\]
Thus $\hat\theta$ is unbiased.
\end{solproblema}

\begin{solproblema}[prob:6f639cf]
\begin{enumerate}[label=(\alph*)]
	\item The theoretical first moment is $E[X]=1/\lambda$. Equating it with $\bar X$ gives $\hat\lambda_{MoM}=1/\bar X$.
	\item The log-likelihood is $\ell(\lambda)=n\ln\lambda-\lambda\sum x_i$. Setting
	\[
		\frac{d\ell}{d\lambda}=\frac n\lambda-\sum x_i=0
	\]
	gives $\hat\lambda_{MLE}=n/\sum x_i=1/\bar X$. The methods coincide here because the one-parameter exponential model is governed by the same sufficient statistic that appears in the first moment equation.
\end{enumerate}
\end{solproblema}

\begin{solproblema}[prob:0363afd]
\begin{enumerate}[label=(\alph*)]
	\item The likelihood is
	\[
		L(p)=p^{\sum x_i}(1-p)^{n-\sum x_i},
	\]
	so
	\[
		\ell(p)=\left(\sum x_i\right)\ln p+\left(n-\sum x_i\right)\ln(1-p).
	\]
	Setting $\ell'(p)=0$ gives
	\[
		\hat p_{MLE}=\frac1n\sum_{i=1}^n X_i=\bar X.
	\]
	\item Since $\Var(X_i)=p(1-p)$,
	\[
		\Var(\hat p_{MLE})=\Var(\bar X)=\frac{p(1-p)}{n}.
	\]
	The estimated standard error is
	\[
		\widehat{\operatorname{SE}}(\hat p_{MLE})=\sqrt{\frac{\hat p(1-\hat p)}{n}}.
	\]
\end{enumerate}
\end{solproblema}

\begin{solproblema}[prob:977b8dd]
\begin{enumerate}[label=(\alph*)]
	\item For one observation,
	\[
		\ln f(X\mid\lambda)=X\ln\lambda-\lambda-\ln(X!).
	\]
	Then
	\[
		\frac{\partial}{\partial\lambda}\ln f=\frac X\lambda-1,\qquad
		\frac{\partial^2}{\partial\lambda^2}\ln f=-\frac X{\lambda^2}.
	\]
	Thus
	\[
		I(\lambda)=-E\left[-\frac X{\lambda^2}\right]=\frac{E[X]}{\lambda^2}=\frac1\lambda.
	\]
	\item For $n$ independent observations, the bound is
	\[
		\operatorname{CRLB}=\frac{1}{nI(\lambda)}=\frac{\lambda}{n}.
	\]
	\item Since $E(\bar X)=\lambda$ and
	\[
		\Var(\bar X)=\frac{\Var(X_i)}n=\frac{\lambda}{n},
	\]
	the estimator $\hat\lambda_{MLE}=\bar X$ is unbiased and attains the CRLB. Therefore it is efficient and UMVUE.
\end{enumerate}
\end{solproblema}

\begin{solproblema}[prob:23bb182]
\begin{enumerate}[label=(\alph*)]
	\item Since $E[X]=\theta/2$, the method-of-moments estimator is $\hat\theta_{MoM}=2\bar X$. The likelihood is $L(\theta)=\theta^{-n}$ for $\theta\ge X_{(n)}$ and $0$ otherwise. Differentiating $\ell(\theta)=-n\ln\theta$ and setting it equal to zero gives no solution because the support depends on $\theta$ and the regularity conditions fail. The likelihood decreases over the admissible region, so $\hat\theta_{MLE}=X_{(n)}$.
	\item The method-of-moments estimator is unbiased. For the maximum,
	\[
		E[X_{(n)}]=\frac{n}{n+1}\theta,
	\]
	so the MLE is biased downward. The unbiased correction is
	\[
		\hat\theta^*=\frac{n+1}{n}X_{(n)}.
	\]
	In practice, $\hat\theta^*$ is attractive because it retains the information in the sample maximum while removing the MLE's systematic downward bias.
\end{enumerate}
\end{solproblema}

\begin{solproblema}[prob:0675127]
For one Bernoulli observation, the log-likelihood is
\[
	\ell(p)=y\ln p+(1-y)\ln(1-p).
\]
The cross-entropy loss is exactly the negative log-likelihood:
\[
	\mathcal L(p)=-\ell(p).
\]
Therefore,
\[
	\arg\min_p \mathcal L(p)=\arg\max_p \ell(p).
\]
Training a binary classifier by minimizing cross-entropy is thus the same maximum-likelihood principle applied to a Bernoulli observation.
\end{solproblema}
